% !TeX root = ../main.tex
%%% ---------------------------------------------------------------------------
%%% 引言
%%%
%%% 结构建议（学术写作通行的「漏斗式」）：
%%%   1. 研究领域的重要性与应用背景
%%%   2. 已有工作及其成果
%%%   3. 已有工作的不足 —— 这是全文的立论点，要具体，不要泛泛说「效果不佳」
%%%   4. 本文的思路与主要贡献（分条列出）
%%%   5. 全文组织
%%%
%%% 引用一律用 \cite{...}；不要手写 [1]。
%%% ---------------------------------------------------------------------------
\section{引言}
\label{sec:intro}

高分辨率遥感图像的目标检测是对地观测、灾害评估与城市管理的关键技术
环节\cite{cheng2016survey}。与自然图像相比，遥感图像的目标具有尺度跨度大、
方向任意、背景杂乱三个显著特点，其中尺度问题尤为突出：同一景中大型建筑
可占据数千像素，而车辆、船舶等目标往往只有十余像素。

围绕这一问题，已有工作主要沿两条路线展开。一是构建多尺度特征表示，
以特征金字塔网络\cite{lin2017fpn}为代表，通过自顶向下的横向连接把深层
语义信息回传到浅层；二是改进检测头与标签分配策略\cite{zhang2020atss}，
使不同尺度的目标获得更合理的正样本。国内学者在此基础上针对遥感场景
提出了旋转框表示与角度回归的改进\cite{wang2021remote}。

然而，现有方法在处理小目标时仍存在两点不足：

\begin{enumerate}
  \item 特征金字塔的横向连接采用固定权重相加，未考虑不同空间位置对各
        层级特征的需求差异——背景区域与小目标区域被同等对待；
  \item 相邻层级之间存在因下采样引入的感受野偏移，直接逐元素相加会
        造成特征错位，这一误差在小目标上被放大。
\end{enumerate}

针对上述问题，本文提出\ourmethod{}。主要贡献如下：

\begin{enumerate}
  \item 提出跨层注意力门控，按空间位置自适应地加权各层级特征，使小目标
        区域更多地保留浅层高分辨率信息；
  \item 引入可变形卷积对齐相邻层级的感受野偏移，缓解逐元素融合造成的
        特征错位；
  \item 在 \dataset{} 上验证了方法的有效性，并通过消融实验分离了两个
        模块各自的贡献。
\end{enumerate}

全文组织如下：\cref{sec:related} 综述相关工作；\cref{sec:method} 给出
\ourmethod{} 的结构与理论分析；\cref{sec:exp} 报告实验结果；
\cref{sec:conclusion} 总结全文。
